% =====================================================================
\chapter{Problemas de optimización y convexidad}
\label{tema:1}
% =====================================================================

Este tema plantea el entrenamiento de un modelo como un problema de minimización sin restricciones y reúne las herramientas de análisis que se usan en el resto del módulo: convexidad, convexidad fuerte, suavidad de Lipschitz y condiciones de optimalidad. Las definiciones y los resultados siguen los capítulos~2, 3 y~9 de~\cite{bv}; la discusión sobre redes profundas sigue~\cite[§8.1 y §8.2]{gbc} y~\cite[§12.1 y §12.2]{d2l}. Parte de la materia de convexidad se introdujo en el módulo~4 (tema~5); aquí se completa con las constantes $\mu$ y $L$, que determinan la velocidad de los algoritmos.

\section{Minimización del riesgo empírico}

\begin{definicion}[Problema de optimización sin restricciones]
Sea $f\colon\R^{n}\to\R$. El problema \emph{minimizar $f(\theta)$, $\theta\in\R^{n}$}, consiste en hallar el \emph{valor óptimo} $f^{\ast}=\inf_{\theta}f(\theta)$ y, si existe, un \emph{minimizador global} $\theta^{\ast}$ con $f(\theta^{\ast})=f^{\ast}$. Un punto $\bar\theta$ es un \emph{minimizador local} si existe $r>0$ tal que $f(\bar\theta)\le f(\theta)$ para todo $\theta$ con $\norm{\theta-\bar\theta}<r$.
\end{definicion}

\begin{definicion}[Riesgo empírico]\label{def:riesgo}
Sean $N$ ejemplos de entrenamiento y $\ell_i(\theta)$ la pérdida del modelo de parámetros $\theta$ sobre el ejemplo $i$. El \emph{riesgo empírico} es
\[
J(\theta)=\frac1N\sum_{i=1}^{N}\ell_i(\theta).
\]
\end{definicion}

En un modelo de lenguaje, cada ejemplo es una posición de un texto y $\ell_i(\theta)=-\ln p_\theta(x_t\mid x_{1:t-1})$ es la log-verosimilitud negativa del token observado; $J$ es la pérdida $\mathcal{J}$ del módulo~1 (tema~5). El objetivo real no es minimizar $J$, sino el riesgo esperado sobre la distribución de los datos, al que $J$ solo aproxima; por eso el aprendizaje se diferencia de la optimización pura~\cite[§8.1]{gbc}. Este módulo trata, sin embargo, del problema de optimización.

\section{Convexidad, convexidad fuerte y suavidad}

\begin{definicion}[Conjunto y función convexos]
Un conjunto $C\subseteq\R^{n}$ es \emph{convexo} si $t x+(1-t)y\in C$ para todo $x,y\in C$ y $t\in[0,1]$. Una función $f\colon\R^{n}\to\R$ es \emph{convexa} si
\[
f\bigl(tx+(1-t)y\bigr)\le t f(x)+(1-t)f(y)\qquad\text{para todo } x,y\in\R^{n},\ t\in[0,1],
\]
y \emph{estrictamente convexa} si la desigualdad es estricta cuando $x\ne y$ y $t\in(0,1)$~\cite[cap.~2 y §3.1]{bv}.
\end{definicion}

\begin{teorema}[Caracterización de primer orden]\label{teo:primer-orden}
Sea $f$ diferenciable. Entonces $f$ es convexa si y solo si
\[
f(y)\ge f(x)+\nabla f(x)\T(y-x)\qquad\text{para todo } x,y\in\R^{n}.
\]
\end{teorema}

Geométricamente, una función convexa queda por encima de todas sus rectas (o hiperplanos) tangentes. Para funciones dos veces diferenciables, la convexidad equivale a $\nabla^{2}f(x)\succeq 0$ para todo $x$ (módulo~4, tema~5; \cite[§3.1]{bv}). Las dos definiciones siguientes refuerzan esta condición por abajo y la limitan por arriba.

\begin{definicion}[Convexidad fuerte]\label{def:fuerte}
Sea $\mu>0$. Una función diferenciable $f$ es \emph{$\mu$-fuertemente convexa} si
\[
f(y)\ge f(x)+\nabla f(x)\T(y-x)+\frac{\mu}{2}\norm{y-x}^{2}\qquad\text{para todo } x,y.
\]
Si $f$ es dos veces diferenciable, esto equivale a $\nabla^{2}f(x)\succeq\mu I$ para todo $x$~\cite[§9.1]{bv}.
\end{definicion}

\begin{definicion}[Suavidad de Lipschitz]\label{def:suave}
Sea $L>0$. Una función diferenciable $f$ es \emph{$L$-suave} si su gradiente es lipschitziano de constante $L$:
\[
\norm{\nabla f(x)-\nabla f(y)}\le L\norm{x-y}\qquad\text{para todo } x,y.
\]
Si $f$ es dos veces diferenciable, esto equivale a que todos los autovalores de $\nabla^{2}f(x)$ estén en $[-L,L]$ para todo $x$; para $f$ convexa, a $\nabla^{2}f(x)\preceq LI$.
\end{definicion}

Si $f$ es $\mu$-fuertemente convexa y $L$-suave, necesariamente $\mu\le L$, y el cociente
\[
\kappa=\frac{L}{\mu}\ge 1
\]
se llama \emph{número de condición} de $f$. En el \cref{tema:2} se demuestra que $\kappa$ controla el número de iteraciones del descenso de gradiente.

\addtocounter{definicion}{1}%

\section{Condiciones de optimalidad}

\begin{teorema}[Condición necesaria de primer orden]\label{teo:fermat}
Si $f$ es diferenciable y $\bar\theta$ es un minimizador local, entonces $\nabla f(\bar\theta)=0$.
\end{teorema}

Un punto con $\nabla f(\bar\theta)=0$ se llama \emph{punto crítico} o \emph{estacionario}. En general, la condición no es suficiente; para funciones convexas, sí lo es.

\begin{teorema}[Optimalidad en el caso convexo]\label{teo:opt-convexa}
Sea $f$ convexa y diferenciable. Un punto $\theta^{\ast}$ es minimizador global si y solo si $\nabla f(\theta^{\ast})=0$. En particular, todo minimizador local de una función convexa es global.
\end{teorema}

\begin{proposicion}[Minimizador de una cuadrática]\label{prop:min-cuadratica}
Sea $f(\theta)=\tfrac12\theta\T Q\theta+c\T\theta+b$ con $Q\succ 0$. Su único minimizador y su valor mínimo son
\[
\theta^{\ast}=-Q^{-1}c,\qquad f^{\ast}=b-\tfrac12\,c\T Q^{-1}c,
\]
y para todo $\theta$ se cumple $f(\theta)-f^{\ast}=\tfrac12(\theta-\theta^{\ast})\T Q(\theta-\theta^{\ast})$.
\end{proposicion}

La identidad de la \cref{prop:min-cuadratica} permite trasladar el minimizador al origen: en lo que sigue, al estudiar cuadráticas se supondrá sin pérdida de generalidad $f(\theta)=\tfrac12\theta\T Q\theta$, con $\theta^{\ast}=0$ y $f^{\ast}=0$.

\begin{proposicion}[Desigualdad de Polyak--Łojasiewicz]\label{prop:pl}
Si $f$ es $\mu$-fuertemente convexa y diferenciable, con valor óptimo $f^{\ast}$, entonces
\[
\norm{\nabla f(x)}^{2}\ge 2\mu\bigl(f(x)-f^{\ast}\bigr)\qquad\text{para todo } x.
\]
\end{proposicion}

La \cref{prop:pl} dice que, lejos del óptimo en valor, el gradiente no puede ser pequeño. Es la propiedad que hace converger linealmente al descenso de gradiente (\cref{teo:gd-fuerte}).

\section{Paisajes de pérdida no convexos}

La pérdida de una red neuronal profunda no es convexa en sus parámetros (módulo~4, tema~5). La información local en un punto crítico la proporciona el hessiano.

\begin{proposicion}[Clasificación de puntos críticos]\label{prop:clasificacion}
Sea $f$ de clase $C^{2}$ y $\bar\theta$ un punto crítico.
\begin{enumerate}
  \item Si $\nabla^{2}f(\bar\theta)\succ 0$, $\bar\theta$ es un minimizador local estricto.
  \item Si $\nabla^{2}f(\bar\theta)$ tiene un autovalor negativo, $\bar\theta$ no es minimizador local; si tiene autovalores de ambos signos, es un \emph{punto de silla}.
  \item Si $\nabla^{2}f(\bar\theta)\succeq 0$ con algún autovalor nulo, el criterio no decide.
\end{enumerate}
\end{proposicion}

La demostración, basada en el desarrollo de Taylor de segundo orden, se dio en el módulo~4 (tema~4). El \cref{tab:criticos} recoge los casos con ejemplos en $\R^{2}$.

\begin{table}[h]
\centering\small
\caption{Puntos críticos en el origen y autovalores del hessiano.}\label{tab:criticos}
\begin{tabular}{@{}llll@{}}
\toprule
Tipo & Autovalores de $\nabla^{2}f$ & Ejemplo & Comportamiento del descenso \\
\midrule
Mínimo local & todos $>0$ & $x^{2}+y^{2}$ & converge si el paso es pequeño \\
Máximo local & todos $<0$ & $-x^{2}-y^{2}$ & se aleja \\
Punto de silla & de ambos signos & $x^{2}-y^{2}$ & se aleja salvo en un conjunto de medida nula \\
Valle plano & alguno nulo & $x^{2}$ & no hay fuerza en la dirección plana \\
\bottomrule
\end{tabular}
\end{table}

\addtocounter{definicion}{1}%

\begin{observacion}[El paisaje de las redes profundas]\label{obs:paisaje}
Se conocen varios hechos empíricos sobre la pérdida de las redes profundas. (i) En alta dimensión, los puntos críticos de pérdida alta son mayoritariamente sillas, no mínimos locales: un mínimo exige que todas las curvaturas sean positivas a la vez~\cite{dauphin}, \cite[§8.2.3]{gbc}. (ii) El espectro del hessiano se concentra cerca de cero, con un grupo menor de autovalores alejados que dependen de los datos~\cite{sagun}; el problema está muy mal condicionado~\cite[§8.2.1]{gbc}. (iii) Los lotes muy grandes tienden a llevar a mínimos más agudos, que se han asociado a una peor generalización~\cite{keskar}. (iv) Las conexiones residuales suavizan visiblemente el paisaje~\cite{li-paisaje}; se estudian en el módulo~11. Los temas siguientes desarrollan algoritmos que funcionan razonablemente en este entorno aunque su análisis se haga en el caso convexo.
\end{observacion}

